\documentclass[a4paper,12pt]{article}
\usepackage{hyperref}
\usepackage{float}
\usepackage{graphicx}
\usepackage{caption}

%% Enable license helpers
\graphicspath{{../cc_beamer/}}
%%%%%%%%%%%%%%%%%%%%%%%%%%%%%%%%%%%%%%%%%%%%%%%%%%%%%%%%%%%%%%%%
%% ccBeamer 0.1, 2007-07-02                                   %%
%% Written by Sebastian Pipping <webmaster@hartwork.org>      %%
%% ---------------------------------------------------------- %%
%% Licensed under Creative Commons Attribution-ShareAlike 3.0 %%
%% http://creativecommons.org/licenses/by-sa/3.0/             %%
%%%%%%%%%%%%%%%%%%%%%%%%%%%%%%%%%%%%%%%%%%%%%%%%%%%%%%%%%%%%%%%%


%% Images
\newcommand{\CcImageBy}[1]{%
	\includegraphics[scale=#1]{creative_commons/cc_by_30.pdf}%
}
\newcommand{\CcImageCc}[1]{%
	\includegraphics[scale=#1]{creative_commons/cc_cc_30.pdf}%
}
\newcommand{\CcImageDevNations}[1]{%
	\includegraphics[scale=#1]{creative_commons/cc_dev_nations_30.pdf}%
}
\newcommand{\CcImageNc}[1]{%
	\includegraphics[scale=#1]{creative_commons/cc_nc_30.pdf}%
}
\newcommand{\CcImageNd}[1]{%
	\includegraphics[scale=#1]{creative_commons/cc_nd_30.pdf}%
}
\newcommand{\CcImagePd}[1]{%
	\includegraphics[scale=#1]{creative_commons/cc_pd_30.pdf}%
}
\newcommand{\CcImageSa}[1]{%
	\includegraphics[scale=#1]{creative_commons/cc_sa_30.pdf}%
}
\newcommand{\CcImageSampling}[1]{%
	\includegraphics[scale=#1]{creative_commons/cc_sampling_30.pdf}%
}
\newcommand{\CcImageSamplingPlus}[1]{%
	\includegraphics[scale=#1]{creative_commons/cc_sampling_plus_30.pdf}%
}


%% Groups
\newcommand{\CcGroupBy}[1]{% zoom
	\CcImageBy{#1}%
}
\newcommand{\CcGroupByNc}[2]{% zoom, gap
	\CcImageBy{#1}\hspace*{#2}\CcImageNc{#1}%
}
\newcommand{\CcGroupByNcNd}[2]{% zoom, gap
	\CcImageBy{#1}\hspace*{#2}\CcImageNc{#1}\hspace*{#2}\CcImageNd{#1}%
}
\newcommand{\CcGroupByNcSa}[2]{% zoom, gap
	\CcImageBy{#1}\hspace*{#2}\CcImageNc{#1}\hspace*{#2}\CcImageSa{#1}%
}
\newcommand{\CcGroupByNd}[2]{% zoom, gap
	\CcImageBy{#1}\hspace*{#2}\CcImageNd{#1}%
}
\newcommand{\CcGroupBySa}[2]{% zoom, gap
	\CcImageBy{#1}\hspace*{#2}\CcImageSa{#1}%
}
\newcommand{\CcGroupDevNations}[1]{% zoom
	\CcImageDevNations{#1}%
}
\newcommand{\CcGroupNcSampling}[2]{% zoom, gap
	\CcImageNc{#1}\hspace*{#2}\CcImageSampling{#1}%
}
\newcommand{\CcGroupPd}[1]{% zoom
	\CcImagePd{#1}%
}
\newcommand{\CcGroupSampling}[1]{% zoom
	\CcImageSampling{#1}%
}
\newcommand{\CcGroupSamplingPlus}[1]{% zoom
	\CcImageSamplingPlus{#1}%
}


%% Text
\newcommand{\CcLongnameBy}{Attribution}
\newcommand{\CcLongnameByNc}{Attribution-NonCommercial}
\newcommand{\CcLongnameByNcNd}{Attribution-NoDerivs}
\newcommand{\CcLongnameByNcSa}{Attribution-NonCommercial-ShareAlike}
\newcommand{\CcLongnameByNd}{Attribution-NoDerivs}
\newcommand{\CcLongnameBySa}{Attribution-ShareAlike}

\newcommand{\CcNote}[1]{% longname
	This work is licensed under the \textit{Creative Commons #1 3.0 License}.%
}


\begin{document}
\title{Reducing iptables configuration complexity using chains}
\author{Dieter Adriaenssens\\
Ghent University, Belgium}
\date{LinuxTag - Berlin\\
May 8th, 2014}
\maketitle

% Creative commons license
\vfill
\begin{center}
  \CcGroupBySa{0.83}{0.95ex}\\[2.5ex]
    {\footnotesize\CcNote{\CcLongnameBySa}}
  \vspace*{-2.5ex}
\end{center}

\pagebreak
\section{Abstract}

After a short introduction of iptables, a local firewall for Linux systems, we go more into detail what chains are, and how these can be used to reduce the complexity of an iptables configuration by grouping similar filtering rules. This results in a better structured configuration that is easier to maintain and that runs faster.\\

The audience should have a basic understanding of network concepts like protocols, ports, IP-addresses and packets.

\tableofcontents
\pagebreak

\section{Introduction}
Setting up a firewall on your *nix box, being it a workstation, laptop, or server, is always a good idea. In most cases, you can do with some simple firewall rules, f.e. on your laptop, block all incoming requests (except the established connections, i.e. the replies on the outgoing requests you made), or on a simple webserver (allow port 80 only).\\

But if you need more complex rules, f.e. a server that hosts a website available for the entire internet, but with an ssh and samba service that should only be available for the local subnet, or even some specific IP addresses, it becomes a bit more complex.\\
And if you want to filter the outgoing traffic as well, your iptables rules get a mess after a while, and when you want to change anything, chances of a mistake or forgetting something are high, which may result in locking yourself out of your box (at least for remote access), or leaving something open that shouldn't.\\

To make your rules more manageable, you can make use of chains in your iptables rules. This paper is inspired by an article that uses chains to make iptables more efficient (faster). The goal of this paper is to get iptables rules that are easier to read and configure, but it will result in faster handling of packets as well.\\

\section{Brief introduction to iptables}
\subsection{ESTABLISHED state}
When using this option, you can filter for established connections. If you define it in both the INPUT and OUTPUT rules, you only have to define in the INPUT rules which NEW incoming requests should be allowed, and in the OUTPUT rules which NEW outgoing request are allowed. The established connections will be allowed and should not be redefined (making the configuration a lot more readable and maintainable). An example allowing only an ssh service without using the ESTABLISHED state would be :
\begin{verbatim}
  iptables -A INPUT -p tcp --dport ssh -j ACCEPT
  iptables -A INPUT -j REJECT
  iptables -A OUTPUT -p tcp --sport ssh -j ACCEPT
  iptables -A OUTPUT -j REJECT
\end{verbatim}

Basically, every incoming/outgoing connection is dropped, except if the incoming packet has port 22 (ssh) as destination, or if the outgoing packet was sent from port 22 (which is the reply of the ssh server).\\

When using ESTABLISHED state, this will become :
\begin{verbatim}
  iptables -A INPUT -p tcp --dport ssh -j ACCEPT
  iptables -A INPUT -j REJECT
  iptables -A OUTPUT -m state --state ESTABLISHED -j ACCEPT
  iptables -A OUTPUT -j REJECT
\end{verbatim}

Now, every incoming/outgoing connection is dropped, except if the incoming packet has port 22 (ssh) as destination, or if the packet belongs to an established connection. Because incoming connections to port 22 are allowed, the firewall will remember a packet coming in, creating a 'connection' for the host/port the packet originates from when the ssh server replies to it. So when the reply of the ssh server is sent out, it matches an 'established' connection and will be allowed out.\\

In this example, the benefit of using the connection state is not clear, but when more allowed incoming services are added, they only have to be added on the INPUT chain, but not on the OUTPUT chain, because they are covered by the ESTABLISHED rule.\\
In the first example (without the ESTABLISHED rule), every allowed incoming connection should be repeated in the OUTPUT chain, matching the packets sent for the outgoing connection, which results in an equal amount of rules on both chains.\\
If you want to do filtering in both directions (allowing incoming request for listening services and outgoing request for remote services), this can become very messy, and almost unmaintainable without making mistakes.\\

\subsection{Introducing chains}
When two services (on different ports) should be available to a limited but identical list of IP addresses.
Without using chains, for every combination of port and IP a rule should be created :
\begin{verbatim}
  iptables -A INPUT -p tcp -m tcp -s 10.100.2.3 --dport 22 -j ACCEPT
  iptables -A INPUT -p tcp -m tcp -s 10.100.2.4 --dport 22 -j ACCEPT
  iptables -A INPUT -p tcp -m tcp -s 10.100.2.7 --dport 22 -j ACCEPT

  iptables -A INPUT -p tcp -m tcp -s 10.100.2.3 --dport 10000 -j ACCEPT
  iptables -A INPUT -p tcp -m tcp -s 10.100.2.4 --dport 10000 -j ACCEPT
  iptables -A INPUT -p tcp -m tcp -s 10.100.2.7 --dport 10000 -j ACCEPT
\end{verbatim}

Resulting in a lot of rules, and when an IP address has to be changed, added or removed, this has to be done for every corresponding rule.\\

When using chains, this can be much easier. Imagine, that you first check if the packet matches the destination port, and if it does, jump to a new chain, where a list of IP addresses is checked. :
\begin{verbatim}
  // create new chain admin_IP
  iptables -N admin_IP

  // add rules to chain admin_IP
  iptables -A admin_IP -s 10.100.2.3 -j ACCEPT
  iptables -A admin_IP -s 10.100.2.4 -j ACCEPT
  iptables -A admin_IP -s 10.100.2.7 -j ACCEPT
  // drop all packets that are not matched by previous rules
  iptables -A admin_IP -j DROP

  // filter ports in INPUT chain
  iptables -A INPUT -p tcp -m tcp --dport 22 -j admin_IP
  iptables -A INPUT -p tcp -m tcp --dport 10000 -j admin_IP
\end{verbatim}

As you can see, there is are several benefits of putting the IP addresses in a separate chain :
\begin{itemize}
    \item the list of IP addresses in the separate chain can be reused for both ports, so they have to be defined only once.
    \item adding/changing/removing an IP address is much easier
    \item there is a better overview of the firewall rules
\end{itemize}

\section{Optimizing configuration with an example setup}
\subsection{Example definition}
\begin{itemize}
  \item A web service should be available from all networks (i.e. internet) on port 80 (http) and 443 (https)
  \item The server can be managed remotely using ssh (port 20) and webmin (port 10000), but only from a limited set of IP addresses (admin PC's).
  \item The server hosts a samba service (several TCP ports), that should only be available from a limited set of IP addresses (admin + webmaster PC's).
  \item Outgoing connections will be filtered, but some services should be allowed (dns, dhcp, smtp, ntp) and some external websites should be available to get updates.
\end{itemize}

\subsection{Actual configuration}

\subsubsection{INPUT chain}

\begin{verbatim}
  iptables -A INPUT -i lo -j ACCEPT
  iptables -A INPUT -m state --state ESTABLISHED -j ACCEPT
  iptables -A INPUT -m state --state RELATED -j ACCEPT
  iptables -A INPUT -p icmp -j icmp_in
  iptables -A INPUT -p tcp -m tcp --dport 22 -j admin_IP
  iptables -A INPUT -p tcp -m tcp --dport 10000 -j admin_IP
  iptables -A INPUT -p tcp -m tcp --dport 80 -j ACCEPT
  iptables -A INPUT -p tcp -m tcp --dport 443 -j ACCEPT
  iptables -A INPUT -p tcp -m tcp --dport 139 -j webmaster_IP
  iptables -A INPUT -p tcp -m tcp --dport 445 -j webmaster_IP
  iptables -A INPUT -j DROP
\end{verbatim}

This is the INPUT chain, allowing :
\begin{itemize}
  \item all local traffic (not leaving the physical PC)
  \item established and related connections
  \item ICMP packets (ping, etc.) are handled in a seperate chain icmp\_in
  \item some services
  \begin{itemize}
    \item SSH (tcp 22) and webmin (tcp 10000) allowed for admins (admin\_IP chain)
    \item website (tcp 80 and 443) for everyone
    \item SMB (tcp 139 and 445) for webmasters (and admins, see definition of webmaster\_IP chain)
  \end{itemize}
  \item everything else is not allowed (dropped)
\end{itemize}

\subsubsection{OUTPUT chain}

\begin{verbatim}
  iptables -A OUTPUT -o lo -j ACCEPT
  iptables -A OUTPUT -m state --state ESTABLISHED -j ACCEPT
  iptables -A OUTPUT -m state --state RELATED -j ACCEPT
  iptables -A OUTPUT -p icmp -m icmp --icmp-type 8 -j ACCEPT
  iptables -A OUTPUT -p udp -m udp --dport 53 -j ACCEPT
  iptables -A OUTPUT -p udp -m udp --dport 67 -j ACCEPT
  iptables -A OUTPUT -p tcp -m tcp --dport 25 -j ACCEPT
  iptables -A OUTPUT -p udp -m udp --dport 123 -j ACCEPT
  iptables -A OUTPUT -p tcp -m tcp --dport 80 -j ext_websites
  iptables -A OUTPUT -j DROP
\end{verbatim}

The OUTPUT chain, allowing :
\begin{itemize}
  \item all local traffic (not leaving the physical PC)
  \item established and related connections
  \item ICMP replies (fe. ping)
  \item some remote services :
  \begin{itemize}
    \item DNS (udp 53) : name service
    \item DHCP (udp 67) : dynamic IP address service
    \item SMTP (tcp 25) : mail service
    \item NTP (udp 123) : time service
  \end{itemize}
  \item external websites (tcp 80), listed in chain ext\_websites
  \item everything else is not allowed (dropped)
\end{itemize}

\subsubsection{icmp in chain}

\begin{verbatim}
  iptables -N icmp_in
  iptables -A icmp_in -p icmp -m icmp --icmp-type 8 -j ACCEPT
  iptables -A icmp_in -p icmp -m icmp --icmp-type 0 -j ACCEPT
  iptables -A icmp_in -p icmp -m icmp --icmp-type 3 -j ACCEPT
  iptables -A icmp_in -p icmp -m icmp --icmp-type 4 -j ACCEPT
  iptables -A icmp_in -p icmp -m icmp --icmp-type 11 -j ACCEPT
  iptables -A icmp_in -p icmp -m icmp --icmp-type 12 -j ACCEPT
  iptables -A icmp_in -j DROP
\end{verbatim}

All allowed incoming ICMP message types.

\subsubsection{admin IP chain}

\begin{verbatim}
  iptables -N admin_IP
  iptables -A admin_IP -s 10.100.2.3 -j ACCEPT
  iptables -A admin_IP -s 10.100.2.4 -j ACCEPT
  iptables -A admin_IP -s 10.100.2.7 -j ACCEPT
  iptables -A admin_IP -j DROP
\end{verbatim}

A list of allowed IP addresses of admin PC's. Everything else is not allowed.

\subsubsection{webmaster IP chain}

\begin{verbatim}
  iptables -N webmaster_IP
  iptables -A webmaster_IP -s 10.100.2.11 -j ACCEPT
  iptables -A webmaster_IP -s 10.100.2.17 -j ACCEPT
  iptables -A webmaster_IP -s 10.100.2.34 -j ACCEPT
  iptables -A webmaster_IP -s 10.100.2.50 -j ACCEPT
  iptables -A webmaster_IP -j admin_IP
\end{verbatim}

A list of allowed IP addresses of webmaster PC's.
At the end of the list, it jumps to the admin\_IP chain, chaining both chains.


\subsubsection{external websites chain}

\begin{verbatim}
  iptables -N ext_websites
  iptables -A ext_websites -d 212.211.132.250 -j ACCEPT
  iptables -A ext_websites -d 212.211.132.32 -j ACCEPT
  iptables -A ext_websites -d 195.20.242.89 -j ACCEPT
  iptables -A ext_websites -d 130.89.149.225 -j ACCEPT
  iptables -A ext_websites -d 86.59.118.153 -j ACCEPT
  iptables -A ext_websites -d 130.89.149.227 -j ACCEPT
  iptables -A ext_websites -d 128.31.0.51 -j ACCEPT
  iptables -A ext_websites -d 86.59.118.153 -j ACCEPT
  iptables -A ext_websites -d 67.228.198.100 -j ACCEPT
  iptables -A ext_websites -d 140.211.166.6 -j ACCEPT
  iptables -A ext_websites -d 140.211.166.21 -j ACCEPT
  iptables -A ext_websites -j LOG
\end{verbatim}

A list of allowed external websites for updates (mirrors of Debian, webmin and Drupal, in this example).
All other requests for external websites are logged. This can be useful for monitoring : notification of abuse, or if you forgot to add an allowed website.

\section{Conclusion}
\begin{itemize}
  \item When checking for the ESTABLISHED state, the number of rules can be reduced
  \item Using chains makes your rules better structured which makes them easier to read and better maintainable
  \item Chains can be reused for several rules
  \item Chains can be chained together
  \item Using chains makes checking rules faster, because it only jumps to a chain when a rule matches. Not all rules have to be checked.
\end{itemize}

\end{document}
